\chapter{Firmware and Software Control Flow}
\label{ch:firmware}

\section{Firmware Architecture and Memory Layout}

The bare-metal firmware is the software side of the hardware--software
co-design and is responsible for every aspect of the benchmark loop:
matrix generation, data preparation, accelerator configuration, timing
measurement, correctness checking, and result reporting. No operating
system is present; the firmware runs directly on PicoRV32 from address
\texttt{0x0000\_0000}. Control begins in \texttt{start.S}, a short
RISC-V assembly bootstrap that sets the stack pointer to the top of the
64\,KB RAM, clears the BSS section, and enters \texttt{main()}.
\texttt{main.c} itself defines the 12 test cases, iterates over them, and
for each case drives matrix generation, the software reference, the
accelerator launch, correctness comparison, and UART reporting. The
accelerator driver \texttt{spmm\_accel.c/.h} wraps the MMIO register
interface behind configuration, start, and polling functions, and a
handful of \texttt{memset}/\texttt{memcpy}/\texttt{memmove} stubs in
\texttt{libc\_stubs.c} cover the minimum set of C library calls required
for bare-metal operation. The linker script maps all sections
(\texttt{.text}, \texttt{.rodata}, \texttt{.data}, \texttt{.bss}) into the
same 64\,KB RAM, with the stack growing downward from the top.

Matrix buffers are allocated statically in the upper portion of the
address space, so that the CSR arrays, dense matrix $B$, and the output
matrices (\texttt{c\_accel} and \texttt{c\_cpu}) occupy contiguous regions
below the stack. The largest benchmark (T6: $64 \times 64 \times 32$,
75\% sparsity, NNZ~$\approx 1024$) consumes roughly
\[
  \underbrace{65 \times 4}_{\texttt{rowptr}} +
  \underbrace{1024 \times 4}_{\texttt{colidx}} +
  \underbrace{1024 \times 4}_{\texttt{values}} +
  \underbrace{64 \times 32 \times 4}_{B} +
  \underbrace{64 \times 32 \times 4}_{C} \approx 29\,\text{KB}
\]
of matrix data, leaving comfortable headroom in the 64\,KB budget for
firmware code and stack.

\section{Runtime Data Preparation}
\label{sec:data-prep}

Dense matrix $B$ is generated from a linear congruential PRNG seeded by
the benchmark case's \texttt{seed} field, with each entry sampled as a
signed integer in the range $[-127, 127]$ (directly usable as INT32 in
the full-precision path, and as the pre-quantisation domain for the INT8
path). Fixing the seed makes $B$ deterministic across runs and across
firmware revisions, which is what makes cycle-count regression comparison
across the three optimisation stages meaningful in the first place.

Sparse matrix $A$ is generated in three phases. Nonzero positions are
placed first, according to the pattern selected by the benchmark case.
For a uniform pattern, each row $i$ receives a target count
$n_i = \lfloor K(1 - s/100) \rfloor$ of nonzeros at column indices
sampled without replacement from the PRNG; for a row-skewed pattern,
roughly 30\% of the rows receive a disproportionate share of the total
NNZ budget, simulating the hub-node behaviour characteristic of social
graphs; and for a clustered pattern, nonzeros are biased toward
rectangular blocks within the matrix, mimicking the blocked sparsity seen
in some finite-element problems. Each placed nonzero then receives a
random signed value in $[-127, 127]$. Finally, the CSR representation is
constructed: column indices within each row are sorted in ascending order
(required for determinism and for correct comparison against the CPU
reference), \texttt{rowptr} is produced by a prefix sum of per-row
nonzero counts, and the sorted indices and values are written to
\texttt{colidx} and \texttt{values}. Sorting column indices also makes
the CSR representation canonical across runs, which simplifies debugging
whenever an accelerator output needs to be compared by hand against the
expected reference.

\section{Software Baseline}

The software reference is deliberately written in the simplest possible
form. The full-precision inner structure is:
\begin{verbatim}
for (row = 0; row < M; row++) {
    for (p = rowptr[row]; p < rowptr[row+1]; p++) {
        col = colidx[p];
        a   = values[p];
        for (n = 0; n < N; n++)
            C[row*N + n] += a * B[col*N + n];
    }
}
\end{verbatim}
\noindent\textbf{Listing~\ref{lst:spmm-cpu}}: PicoRV32 software SpMM
(full precision).
\label{lst:spmm-cpu}

Only the M-extension \texttt{MUL} instruction is used for the inner
multiply. There is no SIMD, no loop unrolling, and no software
prefetching. This simplicity is intentional: the point of the baseline is
not to measure the fastest CPU implementation possible but to give a
credible, readable software reference against which the accelerator's
structural advantage can be isolated. The INT8 CPU reference
(\texttt{spmm\_cpu\_q8}) follows the same structure but operates on
signed 8-bit operands with INT32 accumulation, and it is the version used
for correctness comparison against the INT8 accelerator output.

\section{INT8 Quantisation and Packing}
\label{sec:quant-flow}

Before an INT8 accelerator run, firmware applies symmetric runtime
quantisation separately to the CSR values array and to matrix $B$.
For a vector $\mathbf{v}$ of length $L$, the procedure computes
$v_{\max} = \max_i |v_i|$, sets scale $s = v_{\max}/127$, and quantises
each element as $q_i = \mathrm{clip}(\mathrm{round}(v_i/s), -127, 127)$.
The edge case $v_{\max} = 0$ is handled by setting all quantised values
to zero and $s = 1$. This matches Equation~\ref{eq:quant} from
Section~2.6. The scale factor is computed and then discarded at
quantisation time: since the correctness check compares accelerator
output against the \emph{quantised} CPU reference (which accumulates the
same signed INT8 values), there is no need to store $s$, and keeping the
path scale-free simplifies both the accumulator logic and the verification.

After quantisation, four consecutive INT8 values are packed into a single
32-bit word in little-endian byte order:
\begin{verbatim}
for (i = 0; i < len; i += 4) {
    word = ((uint32_t)(uint8_t)q[i+0])        |
           ((uint32_t)(uint8_t)q[i+1] << 8)   |
           ((uint32_t)(uint8_t)q[i+2] << 16)  |
           ((uint32_t)(uint8_t)q[i+3] << 24);
    dst[i/4] = word;
}
\end{verbatim}

Matrix $B$ is packed row by row, so a single 32-bit RAM read on the
hardware side delivers four aligned INT8 values and directly populates
four lanes of the $B$-segment buffer after sign extension, exactly as
described in Section~\ref{sec:fsm}.

\section{Accelerator Launch and Timing}

The accelerator driver exposes \texttt{accel\_configure()} for writing
dimension and base-address registers, \texttt{accel\_start\_mode()} for
issuing the control write with the \texttt{start} pulse and the
\texttt{relu\_en}/\texttt{int8\_mode} bits, and \texttt{accel\_wait\_done()}
for polling \texttt{STATUS.done}. \texttt{accel\_run\_spmm\_int8()} wraps
these into the canonical launch sequence used on every benchmark case.
Firmware first writes \texttt{clear\_done} to reset any prior completion
state, calls \texttt{accel\_configure()} to set dimensions and base
addresses, captures the pre-launch cycle count $t_0$, issues the
\texttt{start}-mode write, polls \texttt{STATUS.done} in a tight loop,
captures the post-completion cycle count $t_1$, writes \texttt{clear\_done}
once more to reset for the next run, and reports $t_1 - t_0$ as the
accelerator cycle count. The reported figure therefore includes register
write overhead, the full hardware execution time, and the polling loop
delay, and is explicitly an end-to-end offload latency rather than a pure
datapath throughput measurement.

Timing itself relies on the PicoRV32 \texttt{rdcycle} CSR, which returns
the lower 32 bits of a monotonically increasing hardware cycle counter
clocked at 50\,MHz. The read is wrapped in a short inline assembly stub:
\begin{verbatim}
static inline uint32_t read_cycles(void) {
    uint32_t cyc;
    asm volatile("rdcycle %0" : "=r"(cyc));
    return cyc;
}
\end{verbatim}

The 32-bit counter wraps at $2^{32}$ cycles, which at 50\,MHz is
approximately 85 seconds. The largest benchmark in the suite completes
in fewer than two million cycles, so wrap-around is not a concern and no
32-to-64-bit extension is required in firmware.

\section{Reporting and Correctness Verification}

Results are output as CSV-formatted lines over the on-chip UART at
115200\,baud (50\,MHz / 434 cycles per bit), with each line containing
the test ID, the dimensions $M$, $K$, $N$, the sparsity percentage and
pattern, the NNZ count, the CPU and accelerator cycle counts, the derived
speedup, and a PASS/FAIL verdict. A simplified summary is also driven
onto the DE1-SoC's six seven-segment displays (the upper pair shows
CPU cycles scaled by 64, the middle pair shows raw accelerator cycles,
and the lower pair shows integer speedup), so that a running benchmark
can be read at a glance without a host terminal. The LED output mirrors
the busy/done signals, so visually the board is lit during accelerator
execution and cleared the moment computation completes.

Correctness verification is embedded directly in the benchmark loop:
after each accelerator run, firmware performs an element-by-element
comparison between \texttt{c\_accel} and the INT8 CPU reference
\texttt{c\_cpu}:
\begin{verbatim}
for (i = 0; i < M*N; i++) {
    if (c_accel[i] != c_cpu[i]) {
        pass = 0; break;
    }
}
\end{verbatim}
The full-precision CPU result is also computed in parallel but untimed,
and serves as a sanity reference against which quantisation error can be
inspected if needed. Because both the CPU and the accelerator apply the
same symmetric INT8 quantisation to the same input data and accumulate
identical INT32 partial sums, their outputs must match exactly, and any
discrepancy flags a hardware bug rather than a numerical artefact. All
12 cases reported PASS across all three accelerator stages on hardware.
